\documentclass[12pt]{article}
\usepackage[utf8]{inputenc}
\usepackage[T1]{fontenc}
\usepackage{hyperref}
\title{Kuan Language Resources for K-12 Curriculum Development}
\author{Compiled by Claude Sonnet 5 (claude-sonnet-5) for Andr\'as Kornai}
\date{}
\begin{document}
\maketitle
\section{Introduction}
Kuan (ISO 639-3 uan) is a Southwestern Tai (Kra-Dai) language of Laos. Joshua Project places the Kuan in Viangthong District, Borikhamxai Province, near the Mouan River. Wikipedia cites about 886 speakers in the 2015 census, while Joshua Project gives about 1,000. Do not confuse it with the Khuen (Mon-Khmer) or with a Kuan language of southern China.
\section{Primary Educational Resources}
No primers, textbooks or school programmes were found.
\section{Linguistic and Reference Materials}
\begin{itemize}
\item Wikipedia, Kuan language (Laos): \url{https://en.wikipedia.org/wiki/Kuan_language_(Laos)}
\end{itemize}
\section{Religious and Community Resources}
\begin{itemize}
\item Joshua Project, Kuan in Laos (people group profile): \url{https://joshuaproject.net/people_groups/18527}
\item Joshua Project, language profile uan: \url{https://joshuaproject.net/languages/uan}
\end{itemize}
Joshua Project notes that a Bible translation has been started, but it lists no written material.
\section{Digital and Audio Resources}
None found.
\section{Curriculum Development Context}
Documentation is very thin and no teaching materials were found. Any curriculum would have to begin with community-led documentation, and it would probably also need Lao-medium support.
\end{document}